\chapter{Methodology}
\label{chap:methodology}

\section{Introduction}
The primary aim of this research is to design, implement, and evaluate a dynamic, heuristic caching architecture, the Dynamic Traffic-Aware Cache (DTAC), to mitigate the cache pollution caused by Artificial Intelligence (AI) scrapers, thereby preserving the temporal locality and efficiency of web infrastructure for legitimate human users. To achieve this, the study employs an experimental systems research design that bridges software engineering with controlled performance benchmarking.

This chapter outlines the methodology and procedures utilized to construct and validate the proposed system. It is structured to guide the reader through the complete research lifecycle, from initial system design to final metric evaluation. Section 3.2 provides a high-level overview of the research design and system architecture, detailing the data flow between the Heuristic Proxy and the DTAC engine. Section 3.3 outlines the development environment, stating the hardware specifications, software stack, and simulation tools required to ensure the reproducibility of the study.

The core implementation is divided into two distinct phases. Section 3.4 describes Phase 1: System Development, which details the algorithmic design of the passive AI-score generator, the score-aware cache engine that replaces the recency ordering of the standard Least Recently Used (LRU) algorithm, and the calibration of their parameters. Section 3.5 outlines Phase 2: Workload Simulation and Data Generation, explaining the models used to synthesize legitimate Zipfian human browsing sessions and automated scan patterns. Finally, Section 3.6 details the experimental setup and the controlled benchmarking procedures used to test the prototype, while Section 3.7 defines the metrics, such as Cache Hit Ratio (CHR), the Cache Pollution Index (CPI) and computational latency, used to evaluate the DTAC architecture.

\section{Research Design and System Architecture}

\subsection{High-Level Overview}
\begin{figure}[htbp]
    \centering
    \begin{tikzpicture}[
        stage/.style={signal, signal to=east, signal from=west, draw, thick,
            signal pointer angle=140, minimum height=4.2cm, text width=2.3cm,
            align=center, font=\scriptsize, inner xsep=0.5cm}]
        \node[stage, signal from=nowhere] (dev) {
            {\footnotesize\bfseries System\\Development}\\[0.5em]
            Heuristic router\\(log-odds AI-score)\\[0.4em]
            DTAC engine\\(ghost admission,\\score-weighted\\eviction)};
        \node[stage, right=-0.55cm of dev] (gen) {
            {\footnotesize\bfseries Dataset\\Generation}\\[0.5em]
            Zipf human\\sessions and\\uniform scrapers\\(naive, stealthy)\\[0.4em]
            10--75\% bots,\\seeds 0--5};
        \node[stage, right=-0.55cm of gen] (sim) {
            {\footnotesize\bfseries Calibration \&\\Simulation}\\[0.5em]
            Tune on seed 0\\[0.4em]
            Replay seeds 1--5:\\DTAC, ablation,\\LRU, SIEVE,\\S3-FIFO};
        \node[stage, right=-0.55cm of sim] (eval) {
            {\footnotesize\bfseries Performance\\Evaluation}\\[0.5em]
            Hit rates, BHR,\\CPI, router\\accuracy\\[0.4em]
            Latency, memory,\\paired 95\% CIs};
    \end{tikzpicture}
    \caption{Methodology}

    \label{fig:methodology}
\end{figure}

This study employs an experimental systems research design structured around a proxy-cache pipeline. As illustrated in Figure~\ref{fig:methodology}, the methodology is organized into four sequential stages: System Development, Dataset Generation, Calibration and Simulation, and Performance Evaluation. First, the Heuristic Scoring Engine and the DTAC cache engine are built. Next, synthetic web logs are generated by mixing Zipf-distributed human browsing sessions with uniformly distributed AI scraper requests. The tunable constants of both components are then calibrated on a held-out log that is never used for evaluation. During simulation, the evaluation logs are replayed through DTAC, an ablated version of DTAC that ignores the AI-score, a standard LRU baseline, and two modern scan-resistant baselines (SIEVE and S3-FIFO). Finally, the performance evaluation extracts and compares hit rates, cache pollution, latency and memory metrics, with confidence intervals computed over five independently generated workloads.

\subsection{Architectural Components}
\begin{figure}[htbp]
    \centering
    \begin{tikzpicture}[
        font=\footnotesize,
        box/.style={draw, thick, rounded corners=2pt, align=center, inner sep=4pt},
        inner/.style={draw, rounded corners=2pt, align=center, font=\scriptsize, text width=4.6cm, inner sep=3pt},
        flow/.style={-{Stealth[length=2.2mm]}, thick},
        lbl/.style={font=\scriptsize, align=center}]
        % Heuristic Proxy
        \node[inner] (state) {\textbf{State Manager}\\bounded per-IP state (Space-Saving):\\EWMA velocity, pages vs.\ asset fetches};
        \node[inner, below=0.15cm of state] (score) {\textbf{AI-Score Generator}\\$L_0$ + velocity + datacenter ASN\\+ ``Shy Guys'' + missing \texttt{Accept-Language}\\$\rightarrow$ sigmoid $\rightarrow P(\text{bot}) \in [0,1]$};
        \node[box, fit=(state)(score), inner sep=5pt, label={[font=\footnotesize\bfseries]above:Heuristic Proxy}] (router) {};
        % DTAC engine
        \node[inner, below=1.45cm of score] (ghost) {\textbf{Ghost cache} (metadata only)\\score $\mathrel{+}= 1 - P(\text{bot})$, decayed};
        \node[inner, below=0.7cm of ghost] (main) {\textbf{Main cache} (payloads)\\sampled eviction tournament ($N = 8$)};
        \draw[flow] (ghost) -- node[lbl, left] {promote if\\score $\geq 1.5$} (main);
        \node[inner, below=0.15cm of main] (metrics) {\textbf{Metrics Tracker}: hits, misses, latency, evictions};
        \node[box, fit=(ghost)(main)(metrics), inner sep=5pt, label={[font=\footnotesize\bfseries]below:DTAC Engine (lock-sharded)}] (engine) {};
        \draw[flow] (router.south) -- node[lbl, right] {key + $P(\text{bot})$\\(in memory)} (engine.north);
        % Process boundary
        \begin{scope}[on background layer]
            \node[draw, dashed, gray, rounded corners=4pt, fit=(router)(engine), inner xsep=8pt, inner ysep=22pt] (proc) {};
        \end{scope}
        \node[font=\scriptsize, gray, anchor=south west] at (proc.north west) {\texttt{DtacCache} (single process)};
        % External parties
        \node[box, text width=1.9cm, left=1.3cm of proc] (client) {Client /\\simulated\\web logs};
        \node[box, text width=1.9cm, right=1.3cm of proc] (origin) {Origin\\web server};
        \draw[flow] (client.north) |- node[lbl, above, pos=0.75] {HTTP\\request} (router.west);
        \draw[flow] (engine.west) -| node[lbl, below, pos=0.25] {response\\(hit or miss)} (client.south);
        \draw[flow] ([yshift=3pt]engine.east) -| node[lbl, above, pos=0.25] {miss: fetch} ([xshift=-4pt]origin.south);
        \draw[flow] ([xshift=4pt]origin.south) |- node[lbl, below, pos=0.75] {resource} ([yshift=-3pt]engine.east);
    \end{tikzpicture}
    \caption{System Architecture}

    \label{fig:architecture}
\end{figure}

The core of the system relies on a two-stage architectural model positioned between the incoming traffic and the origin web server, as depicted in Figure~\ref{fig:architecture}. Both stages run inside a single process: a \texttt{DtacCache} object owns the router and the cache engine and passes the computed score from one to the other in memory. The system is composed of the following primary components:

\begin{itemize}
    \item \textbf{The Heuristic Proxy (``The Brains''):} This component inspects each incoming request and serves as the analytical engine. It is subdivided into two distinct modules:
    \begin{itemize}
        \item \textit{State Manager:} This module keeps bounded per-IP state across requests: a decayed request-velocity counter and a record of how many pages an IP has fetched and whether it has ever fetched an auxiliary asset (such as \texttt{favicon.ico}, stylesheets or scripts) that a real browser would typically request.
        \item \textit{AI-Score Generator:} Utilizing the data from the State Manager together with the request's own headers and source network, this module calculates a continuous probability score ranging from 0.0 to 1.0 that quantifies how likely the request is to have originated from an automated scraper.
    \end{itemize}
    \item \textbf{The DTAC Cache:} Positioned directly upstream of the origin server, this caching engine receives the requested resource key together with the AI-score. It consists of:
    \begin{itemize}
        \item \textit{DTAC Engine (Score-Aware Cache):} The core memory management engine. Instead of a recency-ordered LRU list, it combines a metadata-only ghost cache for admission, a time-decayed popularity score per resource to which each request contributes in proportion to how human it appears, and a randomized sampling tournament for eviction. Content requested mainly by probable bots therefore accumulates little score, is rarely admitted, and is evicted first.
        \item \textit{Metrics Tracker:} An integrated monitoring component that records hits, misses, per-phase latencies and eviction events, and exports them for the Performance Evaluation phase.
    \end{itemize}
\end{itemize}

\subsection{Data Flow and Request Lifecycle}
The lifecycle of an individual request follows a unidirectional path through the architecture (see Figure~\ref{fig:architecture}):

\begin{enumerate}
    \item \textbf{Ingestion:} An HTTP request, either replayed from the simulated web logs or received by the standalone reverse proxy, arrives at the Heuristic Proxy.
    \item \textbf{Scoring and Forwarding:} The State Manager and AI-Score Generator evaluate the request and compute its 0.0--1.0 bot probability. The score is passed directly, in memory, to the DTAC engine together with the resource key. (The standalone proxy can optionally echo the score back to the client in an \texttt{X-Bot-Probability} header for debugging and benchmarking.)
    \item \textbf{Cache Evaluation:}
    \begin{itemize}
        \item \textit{Cache Hit:} If the DTAC engine already holds the requested resource in its main cache, the resource's score is updated and the cached payload is returned without contacting the origin server.
        \item \textit{Cache Miss:} If the resource is not in the main cache, its metadata is recorded (or updated) in the ghost cache and the request is forwarded to the Origin Web Server.
    \end{itemize}
    \item \textbf{Retrieval and Admission:} The origin returns the resource to the client. The payload is stored in the main cache only once the resource's accumulated, AI-score-weighted popularity crosses a promotion threshold, which in practice means a second request from a probable human. If the main cache is full, a randomized eviction tournament first removes the lowest-scoring resident entry. The origin fetch is performed outside the cache's internal locks, so a slow origin does not block unrelated requests.
\end{enumerate}

\section{Development Environment and Tools}
To ensure the reproducibility of the study, the hardware and software environments were explicitly defined and recorded alongside every experimental run.

\subsection{Hardware Specifications}
All development, testing and benchmarking was carried out on a single consumer laptop featuring a 12th Gen Intel\textsuperscript{\textregistered} Core\textsuperscript{\texttrademark} i5-1235U processor (10 cores: 2 performance and 8 efficiency cores), 16 GB of RAM, and a Solid State Drive (SSD), running Arch Linux.

The same machine was used for the initial coding, functional acceptance testing and algorithm debugging (Phase 1) and for the final workload simulations (Phase 2). To limit the influence of background processes and thermal effects on the measurements, the benchmark runs were executed strictly one at a time rather than in parallel, every reported metric is averaged over five independent workloads, and the comparisons between algorithms are made on identical request sequences (Section 3.6). Cache-efficiency metrics such as hit rates are deterministic functions of the request sequence and the algorithm, and are therefore unaffected by the hardware; only the absolute latency figures depend on it.

\subsection{Software Stack and Programming Language}
The Heuristic Proxy and the Dynamic Traffic-Aware Cache (DTAC) were developed entirely in Rust (stable toolchain, version 1.97). Rust was explicitly selected for this architecture because it bridges the gap between high-level safety and low-level system performance.

\begin{itemize}
    \item \textbf{Memory Safety:} In traditional low-level languages like C and C++, developers have direct control over memory, which allows for exceptional speed but introduces vulnerabilities such as dangling pointers, memory leaks, and buffer overflows. Rust eliminates these bugs at compile time through its strict ``ownership'' and ``borrowing'' rules, guaranteeing memory safety without relying on a runtime garbage collector.
    \item \textbf{High Performance:} Because it compiles directly to machine code and lacks the overhead of garbage collection pauses, Rust achieves low-latency execution speeds comparable to C and C++. This is critical for building a network proxy and cache, where computing the AI-score must not introduce latency that outweighs the time saved by caching.
    \item \textbf{Concurrency:} Rust's type system prevents data races in multi-threaded environments, enabling ``fearless concurrency'' when processing simultaneous HTTP requests.
\end{itemize}

The cache engine is built from the standard library's \texttt{HashMap} and \texttt{Vec} types: each lock shard keeps a hash map from resource key to an index into a dense vector of entries, so that a uniformly random entry can be sampled and removed in constant time (by swapping it with the last element). The State Manager's bounded per-IP table uses a \texttt{HashMap} paired with ordered \texttt{BTreeSet} indices. The asynchronous runtime \texttt{tokio} handles concurrent request processing; the standalone reverse proxy is built on the \texttt{axum} web framework and forwards misses with the \texttt{reqwest} HTTP client. Latencies are recorded into \texttt{hdrhistogram} histograms, and operational metrics are exported through OpenTelemetry to a Prometheus endpoint.

\subsection{Simulation and Log Replay Tools}
Because testing against a live production server introduces unpredictable variables, the workloads were completely simulated. Dedicated tools were developed for data generation, log replay and experiment orchestration:

\begin{itemize}
    \item \textbf{Data Generator (\texttt{generate\_logs.py}):} A Python script (using \texttt{numpy}) that outputs synthetic HTTP request logs in newline-delimited JSON. It blends Zipf-distributed human browsing sessions with uniformly scanning bot traffic at a chosen bot ratio. Each record carries its timestamp, client IP, path, headers, object size and a ground-truth label (human or bot), and the same random seed always produces byte-identical files.
    \item \textbf{Log Replayer (\texttt{replay\_bench}):} A Rust binary that loads a log into memory and replays every record against the chosen cache algorithm \emph{in-process}, with a bounded number of requests in flight concurrently. All time-dependent logic (velocity, score decay) reads a simulated clock that is set to each record's logged timestamp, so the replay can run far faster than real time while the algorithms still observe the log's original timeline. Running in-process removes network variance, so the measured differences reflect cache behaviour rather than the network stack.
    \item \textbf{Networked Validation (\texttt{dtac\_standalone} and \texttt{origin\_stub}):} The replayer can alternatively send every record as a real HTTP request over \texttt{localhost} to the standalone DTAC reverse proxy, which forwards misses to a stub origin server that returns the requested number of bytes. This mode was used to confirm that the in-process results match the behaviour of the full networked system.
    \item \textbf{Experiment Harness (\texttt{run\_experiments.py}):} A Python script that performs the calibration sweeps, generates each seed's logs, runs every cell of the experiment grid sequentially, and aggregates the results into means, confidence intervals and paired differences.
\end{itemize}

\section{System Implementation}
This section details the algorithmic design and software engineering required to translate theoretical bot detection into a high-performance, deployable network component. The system is divided into two interconnected components: a Heuristic Proxy that evaluates incoming traffic via bounded-memory Log-Odds classification, and a Dynamic Traffic-Aware Cache (DTAC) that leverages randomized eviction, a metadata-only ghost cache, and resource-specific time-decayed scoring to limit automated cache pollution. Every per-request operation in both components has a small, bounded cost that does not grow with the volume of traffic observed: the router's state updates take $\mathcal{O}(\log K)$ time for a fixed tracking capacity $K$, and the cache engine's lookups take constant time with evictions costing $\mathcal{O}(N)$ for a small, fixed sample size $N$.

\subsection{The Heuristic Proxy and AI-Score Generator}
The Heuristic Proxy acts as the initial ingress point for all network traffic. Its primary function is to classify incoming HTTP requests probabilistically without relying on heavy client-side challenges. To keep the scoring cost to microseconds, this was implemented through a Composite Trust Scoring architecture relying on a continuous Log-Odds model rather than rigid binary thresholds \cite{chandra_composite_2026}.

\textbf{The Bounded-Memory State Manager:}
To evaluate behavioral anomalies across a session without succumbing to memory exhaustion from internet-scale traffic, the proxy does not use an unbounded \texttt{HashMap} keyed by IP address. Instead, the State Manager implements the Space-Saving algorithm \cite{metwally_efficient_2005} to address the ``Top-K Elephant Flow'' problem commonly encountered in network measurement (Gong et al., 2018). By utilizing a Rust \texttt{HashMap} paired with ordered \texttt{BTreeSet} indices, the proxy maintains strict bounds on memory ($\mathcal{O}(K)$ space for a fixed capacity of $K$ tracked IPs, with $\mathcal{O}(\log K)$ updates). When the table is full, the IP with the smallest count is replaced and its count is inherited by the newcomer, as the algorithm prescribes. Counts decay with a configurable half-life (one hour), so the structure tends to retain persistent high-volume clients (``elephant flows'') while low-volume visitors (``mice flows'') are naturally displaced.

Within this bounded state, the proxy tracks request velocity using an Exponentially Weighted Moving Average (EWMA). Because human request inter-arrival times naturally follow heavy-tailed distributions characterized by micro-bursts \cite{li_queuing_2018}, the proxy maintains an instantaneous velocity score $v$ using a time-decayed counter rather than storing arrays of timestamps. The counter is a fixed-point integer scaled by 1024; on each request it is halved, via a base-2 bitwise right shift, once for every elapsed velocity tick (5 seconds after calibration) and then incremented by one request. After ten idle ticks the counter is reset. This accommodates natural human burstiness while penalizing sustained automated scraping at a constant cost per request.

\textbf{The AI-Score Generation (Log-Odds Classification):}
Instead of a rudimentary weighted sum, the proxy calculates an aggregate score in the Log-Odds domain, in the style of a Naive Bayes classifier. This allows the proxy to begin with a baseline prior probability of automated traffic and perform efficient additive arithmetic. The prior log-odds $L_0$ is derived from an explicit assumed bot base rate $p_0$:

\begin{equation}
L_0 = \ln\left(\frac{p_0}{1 - p_0}\right)
\end{equation}

With the calibrated $p_0 = 0.02$, $L_0 \approx -3.89$: every request is assumed human until the evidence says otherwise, and a clean human request scores a bot probability of about 0.02.

For every triggered heuristic $i$, a weight $W_i$ is added to the running log-odds score. In an idealized Naive Bayes model this weight would be the log-likelihood ratio of observing the evidence $E_i$ from a bot versus a human:

\begin{equation}
W_i = \ln\left(\frac{P(E_i \mid \text{Bot})}{P(E_i \mid \text{Human})}\right)
\label{eq:llr}
\end{equation}

Estimating these ratios requires labelled real-world traffic, which was not available to this study. Estimating them from the synthetic workload instead would be circular, because the workload generator was written to trigger exactly these heuristics (Section 3.4.3). The weights used were therefore \emph{calibrated} on a held-out workload to maximize the separation between human and bot scores under a false-positive constraint, and Equation~\ref{eq:llr} serves as the motivation for the additive model rather than as the method for obtaining its weights.

The heuristic evaluation engine processes four independent signals:
\begin{itemize}
    \item \textbf{Passive Structural Anomalies (``Shy Guys''):} Following \citeA{van_boxem_shy_2026}, the proxy penalizes IPs that fetch pages but never fetch auxiliary assets (favicon, stylesheets, scripts, images or fonts). Browsers request such assets once and then cache them, so the signal is that an IP has \emph{never} fetched one, not that it skipped one on its last page. A grace period of 12 page requests prevents short human visits from being flagged. The heuristic adds $W_{fav} = 2$.
    \item \textbf{Header Anomalies:} Requests that lack an \texttt{Accept-Language} header, which every mainstream browser sends, receive $W_{hdr} = 4$. Other header checks (such as \texttt{Sec-CH-UA} consistency or header ordering) were not implemented.
    \item \textbf{Infrastructure Origin:} The client IP is looked up in an offline copy of the public iptoasn \texttt{ip2asn} database, which maps IP ranges to Autonomous System Numbers (ASNs), using binary search over the sorted ranges ($\mathcal{O}(\log n)$). If the ASN belongs to a curated list of cloud and hosting providers, the request receives $W_{asn} = 3$, reflecting the established use of datacenter origin as a baseline filter in automated traffic detection \cite{martinez_llamas_balancing_2025}.
    \item \textbf{Graded Request Velocity:} Rather than applying a binary step-function \cite{duskin_distinguishing_2009}, the tracked EWMA velocity $v$ applies a continuously scaling penalty inspired by the Weber-Fechner Law. Once a baseline human burst threshold $v_{\tau}$ is exceeded, a graded log-weight is applied:
\end{itemize}

\begin{equation}
W_{velocity} = \max\left(0, \alpha \cdot \ln\left(\frac{v}{v_{\tau}}\right)\right)
\end{equation}

with the calibrated $v_{\tau} = 3$ (in units of the decayed per-tick counter) and $\alpha = 2$.

Once all heuristic weights are summed ($L = L_0 + \sum W_i$), the final log-odds score is converted back to a continuous probability $P(\text{bot})$ ranging from 0.0 to 1.0 using the sigmoid function:

\begin{equation}
P(\text{bot}) = \frac{1}{1 + e^{-L}}
\end{equation}

This value expresses the router's confidence that the request originated from an automated scraper and is passed directly to the DTAC engine. The router also reports which heuristics fired for each request, which allows the evaluation to measure how often each signal fires on human and on bot traffic.

Algorithm 1 illustrates the complete ingress pipeline, demonstrating how the bounded-memory Space-Saving structure and Log-Odds heuristic scoring operate for each incoming request.

\begin{verbatim}
Algorithm 1: Ingress Proxy Pipeline & Log-Odds AI-Score Generation
-------------------------------------------------------------------
Input: incoming_request
Global State:
    SpaceSaver: Bounded per-IP tracker (HashMap + BTreeSet)
    L_0: Baseline prior log-odds, ln(0.02/0.98) = -3.89
    W_asn = 3, W_fav = 2, W_hdr = 4: calibrated weights
    TICK = 5 s, BURST_THRESHOLD = 3, ALPHA = 2, SHY_PAGES = 12

1:  client_ip <- incoming_request.IP
2:  now <- clock.now()
3:
4:  // --- PHASE 1: State Retrieval & Bounded Memory Management ---
5:  // O(log K); evicts the lowest-count IP when the table is full
6:  ip_state <- SpaceSaver.get_or_add_and_evict(client_ip)
7:
8:  // --- PHASE 2: EWMA Velocity Calculation (Bitwise) ---
9:  ticks <- (now - ip_state.last_seen) / TICK
10: if ticks > 10 then
11:     new_velocity <- 1024  // Reset, plus 1 request (x1024)
12: else
13:     decayed_v <- ip_state.velocity >> ticks  // Bitwise shift
14:     new_velocity <- decayed_v + 1024
15: end if
16:
17: // Update structural history and commit state
18: if is_auxiliary_asset(incoming_request.path) then
19:     ip_state.fetched_asset <- TRUE
20: else
21:     ip_state.pages <- ip_state.pages + 1
22: end if
23: ip_state.last_seen <- now
24: ip_state.velocity <- new_velocity
25: SpaceSaver.update(client_ip, ip_state)
26:
27: // --- PHASE 3: Log-Odds Heuristic Evaluation ---
28: score <- L_0
29:
30: // Heuristic A: Continuous Graded Velocity Penalty
31: float_v <- new_velocity / 1024.0
32: if float_v > BURST_THRESHOLD then
33:     score <- score + (ALPHA * ln(float_v / BURST_THRESHOLD))
34: end if
35:
36: // Heuristic B: Infrastructure Origin (Datacenter ASN)
37: if ASN_Database.is_datacenter(client_ip) then
38:     score <- score + W_asn
39: end if
40:
41: // Heuristic C: "Shy Guys" Passive Behavioral Check
42: if ip_state.pages >= SHY_PAGES
        and not ip_state.fetched_asset then
43:     score <- score + W_fav
44: end if
45:
46: // Heuristic D: Missing Accept-Language Header
47: if incoming_request.missing_header("Accept-Language") then
48:     score <- score + W_hdr
49: end if
50:
51: // --- PHASE 4: Sigmoid Conversion & Forwarding ---
52: bot_probability <- 1.0 / (1.0 + exp(-score))
53:
54: return DTAC_Engine.access(incoming_request.path,
                           bot_probability)
\end{verbatim}

\subsection{The Dynamic Traffic-Aware Cache (DTAC) Engine}
The second component of the system is the DTAC engine, positioned upstream of the origin server. Rather than utilizing deterministic eviction algorithms (such as LRU or LFU), which either fail to account for long-term popularity or suffer from sequential scan pollution under bot traffic \cite{zhang_rethinking_2026}, the DTAC engine was implemented using a Randomized Eviction architecture combined with a resource-specific decayed popularity score and a metadata-only Ghost Cache.

\textbf{Resource-Based Decayed Score \& Lock Sharding:}
Unlike the State Manager (Section 3.4.1), which tracks the behavioral velocity of client IP addresses, the DTAC engine tracks the popularity score of the requested resources (e.g., specific URLs or assets).

Whenever a resource is requested, whether it is in the main cache or the ghost cache, its score $s$ is first decayed for the elapsed time and then incremented by $1 - P(\text{bot})$, the complement of the request's AI-score. A request from a confident human (bot probability $\approx 0.02$) adds almost a full point, whereas a highly suspicious request adds effectively zero. This increment is the mechanism that connects the router's output to caching behavior. Decay is applied lazily, only when an entry is touched, in discrete half-life steps:

\begin{equation}
s(t) = s(t_0) \cdot 2^{-\lfloor (t - t_0) / h \rfloor}
\end{equation}

where $t_0$ is the entry's last access and $h$ is the decay half-life. After 32 half-lives the score is treated as zero. With the calibrated half-life of $h = 4$ hours, which exceeds the duration of the benchmark logs, the score behaves much like a human-weighted frequency count, so the policy approaches a Least Frequently Used (LFU) ordering with recency as a tie-breaker. A shorter half-life would make the cache adapt faster when popularity shifts, at the cost of forgetting stable popularity.

Updating this resource score on every read introduces a known concurrency bottleneck: the ``read-path write penalty.'' To prevent concurrent threads from blocking one another while updating the metadata of a popular asset, the system implements lock sharding. Incoming requests are hashed by resource key to one of several independent, mutex-protected shards, each with its own share of the capacity, its own ghost cache and its own eviction. The shard count is chosen automatically as the largest power of two up to 256 that leaves at least 64 entries per shard (for example, 128 shards at a capacity of 10\,000 entries); very small shards would distort the sampling-based eviction described below.

\textbf{Cold-Start Mitigation via Metadata-Only Ghost Cache:}
To resolve the ``Cold Start'' problem, where brand new items are evicted before they can establish a popularity score, while conserving memory, DTAC implements a two-stage admission design inspired by the S3-FIFO caching model \cite{yang_fifo_2023}.

When a cache miss occurs, the system does not store the full data payload. Instead, it records only the resource's identifier, last access time and initial score inside the Ghost Cache (a lightweight, metadata-only structure), and the request is served from the origin. If the resource receives subsequent requests and its decayed score crosses a promotion threshold (calibrated to 1.5), it graduates: the payload is fetched from the origin and stored in the Main Cache. Because a clean human request contributes about 0.98, two human requests promote an item, whereas a scan of requests that each score near 1.0 never accumulates enough score to be admitted. One-off scraper requests therefore never consume main-cache memory.

The ghost cache is itself bounded, at eight times the main-cache capacity per shard. When it is full, a ghost entry is evicted with the same score-aware tournament as the main cache. Entries created by bot requests start with a score near zero and therefore lose first, so a uniform bot scan cannot push popular human items out of the ghost before their second request, as a plain FIFO ghost queue would allow.

The origin fetch for a promotion is performed after the shard lock has been released. Once the payload arrives, the lock is re-acquired and the engine re-checks whether another concurrent request has already promoted the same resource, so a slow origin never blocks other requests and a popular item is never inserted twice.

\textbf{The Eviction Tournament (The Power of N Choices):}
When the Main Cache reaches capacity, the system triggers a randomized eviction tournament instead of maintaining a globally sorted priority queue. The algorithm samples $N$ random entries from the shard (calibrated to $N = 8$), lazily calculates their currently decayed scores, and evicts the entry with the lowest score, breaking ties in favor of evicting the entry with the oldest last access. Once scores have decayed to zero, eviction therefore falls back to LRU order instead of being random. This approach mirrors modern sampling-based eviction frameworks such as Google's HALP (Google Research, 2023) and RankCache (Beckmann et al., 2018). While sampling introduces a minor statistical variance compared to sorting the entire cache, each eviction costs only $\mathcal{O}(N)$ for a small, fixed $N$, independent of cache size, and requires no global ordering to be maintained on every hit.

\begin{verbatim}
Algorithm 2: DTAC Admission, Eviction, Sharding and Ghost Cache
-------------------------------------------------------------------
Input: requested_resource, request_bot_prob (Float 0.0 to 1.0)
Parameters: N_SAMPLES = 8, PROMOTION_THRESHOLD = 1.5,
            HALF_LIFE = 4 h, GHOST_CAP = 8 x shard capacity
State per shard: MainCache (payload + metadata),
                 GhostCache (metadata only)

1:  // Phase 1: Concurrency Management (Lock Sharding)
2:  shard <- shards[hash(requested_resource) % NUM_SHARDS]
3:  acquire_lock(shard)
4:  increment <- 1.0 - request_bot_prob
5:
6:  // Phase 2: Main Cache Hit
7:  if shard.MainCache.contains(requested_resource) then
8:      meta <- shard.MainCache.metadata(requested_resource)
9:      meta.score <- decay(meta.score, now - meta.last_access)
                      + increment
10:     meta.last_access <- now
11:     payload <- shard.MainCache.payload(requested_resource)
12:     release_lock(shard)
13:     return HIT(payload)
14: end if
15:
16: // Phase 3: Miss -- update or create the ghost entry
17: if shard.GhostCache.contains(requested_resource) then
18:     g <- shard.GhostCache.metadata(requested_resource)
19:     g.score <- decay(g.score, now - g.last_access) + increment
20:     g.last_access <- now
21: else
22:     if shard.GhostCache.size() >= GHOST_CAP then
23:         evict_by_tournament(shard.GhostCache)
24:     end if
25:     g <- shard.GhostCache.insert(requested_resource, now,
                                     increment)
26: end if
27: promote <- (g.score >= PROMOTION_THRESHOLD)
28: release_lock(shard)
29:
30: // Phase 4: Origin fetch outside the lock
31: payload <- fetch_from_origin(requested_resource)
32: if promote then
33:     acquire_lock(shard)
34:     if not shard.MainCache.contains(requested_resource) then
35:         if shard.MainCache.is_full() then
36:             evict_by_tournament(shard.MainCache)
37:         end if
38:         g <- shard.GhostCache.remove(requested_resource)
39:         shard.MainCache.insert(requested_resource, payload, g)
40:     end if  // else: a concurrent request already promoted it
41:     release_lock(shard)
42: end if
43: return MISS(payload)
44:
45: // --- Helper: stepwise half-life decay ---
46: function decay(score, delta_t):
47:     k <- floor(delta_t / HALF_LIFE)
48:     if k >= 32 then return 0
49:     return score / 2^k
50:
51: // --- Helper: Eviction Tournament ---
52: function evict_by_tournament(store):
53:     candidates <- store.random_samples(N_SAMPLES)
54:     victim <- argmin over candidates of
55:           (decay(c.score, now - c.last_access), c.last_access)
56:     store.swap_remove(victim)  // O(1) removal from dense vector
\end{verbatim}

\subsection{Tuning Parameters and System Calibration}
The efficacy of both the Heuristic Proxy and the DTAC engine relies heavily on the calibration of their numerical thresholds. Because internet traffic follows heavy-tailed Zipfian distributions \cite{breslau_web_1999}, poorly chosen values lead either to excessive false positives (penalizing humans) or to cache pollution. The system's parameters must therefore be tuned empirically against the target network's traffic. In this study every constant is exposed through JSON configuration rather than hard-coded, and all of them were calibrated on a dedicated tuning workload (seed 0) that is never used in the reported evaluation (seeds 1--5).

The configuration is governed by three primary categories.

\textbf{1. Velocity and Bounded-Memory Thresholds (State Manager and Engine)}
The parameters governing the EWMA velocity counter, the score decay, and the ghost cache dictate the proxy's sensitivity to automated scrapers and the engine's memory of popularity.
\begin{itemize}
    \item \textbf{Velocity Tick and Human Burst Threshold ($v_{\tau}$):} The velocity counter is halved once per tick, and the threshold is the decayed counter value below which no penalty is applied. Initial values of a 1-second tick and $v_{\tau} = 10$ never flagged the naive bots in the benchmark (whose counters settled near 7). After calibration, a 5-second tick with $v_{\tau} = 3$ flags nearly all naive bots, at the cost of the velocity heuristic also firing on about 4\% of human requests from fast multi-page sessions, never enough on its own to push a human above a bot probability of 0.5.
    \item \textbf{Velocity Scaling Factor ($\alpha$):} Determines the aggressiveness of the continuous Log-Odds penalty once an IP exceeds $v_{\tau}$. It was kept at $\alpha = 2$ and not swept.
    \item \textbf{Score Decay Half-Life ($h$):} The initial half-life of one second made scores vanish between two page views of the same session, so the main cache stayed almost empty (25--130 of its 10\,000 slots). Human hit rate rose with every increase of the half-life up to the largest value tested, 4 hours.
    \item \textbf{Ghost Cache Promotion Threshold:} The decayed score a resource must accumulate in the metadata-only Ghost Cache before its payload is admitted. The calibrated value of $1.5$ means ``promote on the second human request''. A threshold of $0.5$ (admit on the first human request) performed worse at a capacity of 10\,000 entries because one-off human requests then competed with popular items.
    \item \textbf{Ghost Cache Capacity:} Eight times the main-cache capacity; any multiplier of 2 or more improved human hit rate by about 1.5 points over a ghost the same size as the main cache.
\end{itemize}

\textbf{2. Probabilistic Heuristic Weights (Log-Odds Engine)}
\begin{itemize}
    \item \textbf{Prior Probability ($L_0$):} The baseline assumption of automated traffic, expressed through the assumed bot base rate $p_0$. For a heavily targeted public API where 60\% of traffic is automated, the prior log-odds would be $L_0 = \ln(0.60 / 0.40) = +0.405$. Calibration favored the low prior $p_0 = 0.02$ ($L_0 \approx -3.89$), because it lets a clean human request contribute almost a full point ($0.98$) to a resource's score. Higher priors (0.3) pushed returning human visitors above a bot probability of 0.5 once a single heuristic fired.
    \item \textbf{Evidence Weights ($W_i$):} As discussed in Section 3.4.1, these were chosen to maximize the gap between the mean bot probability of bots and that of humans, subject to keeping the false-positive rate for returning human visitors at or below 5\%. The calibrated weights ($W_{hdr} = 4$, $W_{asn} = 3$, $W_{fav} = 2$) are much smaller than the empirical log-likelihood ratios measured on the synthetic data (up to $+12$), because a sum of such large terms would saturate every bot at 1.0 and the weights would stop mattering. If a network naturally serves a high volume of enterprise VPN users, the weight for the datacenter heuristic should be decreased to limit false positives.
\end{itemize}

\textbf{3. Eviction Tournament Sample Size ($N$)}
An important parameter in the DTAC engine is $N$, the number of random items sampled during an eviction tournament. If $N$ is too small, the system risks evicting a highly popular item. If $N$ is too large, the engine wastes CPU cycles evaluating decayed scores, defeating the purpose of randomized eviction.

To reason about a suitable $N$, the framework of Order Statistics can be used to evaluate the expected value of the minimum, $E[X_{(1)}]$. In a cache, the distribution of item scores is typically bimodal: a large cluster of low-scoring ``junk'' items (decayed or one-off requests) and a smaller cluster of high-scoring popular items. This can be modeled as a Gaussian Mixture Model (GMM).

While there is no clean algebraic formula for the exact expected minimum of a GMM, the ``sweet spot'' for $N$ can be determined logically. As $N$ increases, the expected minimum drops sharply at first but then hits a point of extreme diminishing returns. This occurs because once the algorithm draws at least one sample from the ``lower mode'' (the junk items), the minimum is effectively bounded by that lower distribution.

Therefore, selecting $N$ is a function of confidence. Let $p$ be the proportion of the cache population that belongs to the lower distribution (evictable items). The probability of failing to pick from this lower distribution in $n$ tries is $(1-p)^n$. The probability $P$ of successfully sampling the lower distribution at least once is:

\begin{equation}
P = 1 - (1-p)^n
\end{equation}

To find an $N$ that achieves a specific confidence level, we solve for $n$:

\begin{equation}
n = \frac{\ln(1 - P)}{\ln(1 - p)}
\end{equation}

\textit{Estimating $N$:} Suppose the cache is heavily utilized, and only 20\% of the cache consists of evictable junk items ($p = 0.20$). To be 95\% confident ($P = 0.95$) that the random sample includes at least one of these candidates, the required sample size is:

\begin{equation}
n = \frac{\ln(1 - 0.95)}{\ln(1 - 0.20)} = \frac{\ln(0.05)}{\ln(0.80)} = \frac{-2.995}{-0.223} \approx 13.4
\end{equation}

Under this assumption, $N = 14$ would give a 95\% probability of finding a highly decayed item without the cost of sorting the entire cache.

For environments where the cache churn is higher and junk items make up a larger proportion of the cache (e.g., $p = 0.50$), the required sample size drops drastically. Achieving 95\% confidence in this scenario requires only $N = 5$ ($1 - 0.5^5 = 0.968$). This aligns with the default parameters established by industry implementations; Redis utilizes $N = 5$ for its \texttt{allkeys-lfu} policy as a balance between CPU efficiency and eviction accuracy (Redis Core Architecture).

The analysis therefore bounds the useful range of $N$ to roughly 5--14. Empirical calibration within this range selected $N = 8$: it was consistently better than $N = 5$, while $N = 14$ was within 0.03 percentage points of it on the tuning workload, confirming the diminishing returns predicted above.

\textbf{Calibration Procedure:}
The constants were calibrated by coordinate descent on the tuning seed, on the naive workload (Section 3.5) at 30\% and 75\% bot traffic and a cache capacity of 10\,000 entries. Each group of parameters was swept on a coarse grid with all others held fixed, followed by a refinement around the chosen value; a new value replaced the current one only if it improved the objective by more than a small tolerance. The router was calibrated first, maximizing the class-mean gap (the mean bot probability of bot requests minus that of human requests) subject to a returning-visitor false-positive rate of at most 5\%. The engine was then calibrated with the router fixed, maximizing human hit rate. Table~\ref{tab:calibration} lists the resulting values. The constants were not re-tuned for other cache capacities; a check at the other capacities found that at 20\,000 entries a promotion threshold of 0.5 would have performed about 2.3 percentage points better than the calibrated 1.5.

\begin{table}[htbp]
    \centering
    \caption{Initial and calibrated values of the router and engine constants}
    \label{tab:calibration}
    \begin{tabular}{llll}
        \toprule
        Parameter & Component & Initial & Calibrated \\
        \midrule
        Prior bot rate $p_0$ & Router & 0.1 & 0.02 \\
        $W_{asn}$ / $W_{fav}$ / $W_{hdr}$ & Router & 1.5 / 1.2 / 1.0 & 3 / 2 / 4 \\
        Velocity tick / $v_{\tau}$ & Router & 1 s / 10 & 5 s / 3 \\
        Shy Guys page threshold & Router & 5 & 12 \\
        Score half-life $h$ & Engine & 1 s & 4 h \\
        Promotion threshold & Engine & 2.0 & 1.5 \\
        Sample size $N$ & Engine & 5 & 8 \\
        Ghost capacity multiplier & Engine & 4 & 8 \\
        \bottomrule
    \end{tabular}
\end{table}

\section{Workload Simulation and Data Generation}
Testing caching algorithms against a live production server introduces uncontrolled variables (such as network latency spikes, hardware throttling, and unpredictable user behavior) that make precise metric evaluation nearly impossible, and labelled production traffic with ground truth about which requests come from bots was not available. To ensure that the performance of the DTAC algorithm was accurately and reproducibly measured, synthetic workloads were generated with the Python data generator (Section 3.3.3). Every log is fully determined by its random seed; five evaluation seeds (1--5) and one tuning seed (0) were generated.

\subsection{Simulating Legitimate Human Traffic}
Human web browsing naturally exhibits temporal locality and highly skewed popularity patterns, heavily favoring a small subset of ``hot'' web pages. To simulate this behavior, the data generator draws the pages that humans request from a Zipfian distribution \cite{breslau_web_1999}.

In a Zipfian distribution, the probability $P(k)$ of a request for the $k$-th most popular web object out of a total catalog size of $N$ items is defined as:

\begin{equation}
P(k) = \frac{1/k^\alpha}{\sum_{n=1}^{N} (1/n^\alpha)}
\end{equation}

For the simulation, the catalog size ($N$) was established at 100,000 unique web objects. The skewness parameter ($\alpha$), which dictates the intensity of the popularity concentration, was set to 0.8, a standard baseline for mirroring realistic human web traffic. Every object has a fixed size drawn once from a lognormal distribution with a median of 16 KB, clipped to the range 512 B--2 MB, so byte-level metrics can be computed.

Human traffic is generated as browser \emph{sessions} rather than as independent requests. Sessions arrive as a Poisson process; each session uses a single IP address, views a geometrically distributed number of pages (mean 5) drawn from the Zipf catalog, and waits an exponentially distributed think time (mean 8 seconds) between pages. Session arrivals are frequent enough to keep the aggregate human rate at about 12.5 requests per second. Every request carries standard browser headers, including a browser \texttt{User-Agent} and \texttt{Accept-Language}. Seventy percent of sessions are \emph{new visitors}, who fetch \texttt{/favicon.ico} once after their first page; the remaining 30\% are \emph{returning visitors}, who already have the favicon cached and never fetch it. Returning visitors who view 12 or more pages are therefore realistic false positives for the ``Shy Guys'' heuristic, which allows the evaluation to measure its cost to humans. Human IP addresses never fall in datacenter ranges.

Crucially, for a given seed the generator produces a single human request stream that is reused, unchanged, in every bot-ratio file and in both bot workloads. Differences in human hit rate across bot ratios and across algorithms are therefore caused by the bots and the algorithms alone.

\subsection{Simulating AI Scraper Traffic}
Conversely, automated AI scrapers and LLM data harvesters utilize ``scan'' patterns characterized by high-volume, single-use access over broad swaths of a website \cite{zhang_rethinking_2026}. To simulate this, bot requests are drawn from a uniform distribution over the \emph{entire} catalog, overlapping with the pages humans request. In a uniform distribution, the probability of requesting any specific item is equal, resulting in very little temporal recurrence and effectively mimicking the one-hit-wonder behavior of scrapers. Two bot workloads were generated.

\textbf{Naive scrapers} are imbued with the mechanical signatures the Heuristic Proxy is designed to detect:
\begin{itemize}
    \item They never request auxiliary assets such as \texttt{favicon.ico}.
    \item They omit browser headers, including \texttt{Accept-Language} and \texttt{Sec-CH-UA}.
    \item Each bot IP sends at a steady, mechanical rate of 3.6 requests per second ($\pm$5\% jitter), in 30-second on/off cycles. The pool of IPs grows with the bot ratio, so that each IP's rate is the same at every ratio and only the number of IPs changes.
    \item About 64\% of bot IPs fall in IP ranges that the ASN fixture used for the benchmark maps to datacenter providers; the rest use reserved documentation (TEST-NET) addresses that match no ASN.
\end{itemize}

The ASN fixture shipped with the implementation was written to cover the generator's datacenter bot ranges, so the datacenter signal is correct by construction in this benchmark. It demonstrates that the lookup works and how much a correct ASN signal helps the cache, not that ASN detection works on real traffic.

\textbf{Stealthy scrapers} imitate browser sessions in order to evaluate where the heuristics stop working. They send full browser headers (including \texttt{Accept-Language}, \texttt{Sec-CH-UA} and a human \texttt{User-Agent}) from fresh, residential-looking IP addresses, issue only about 5 requests per IP with exponential think times, and fetch \texttt{/favicon.ico} once per IP, which defeats the ``Shy Guys'' heuristic. They still scan the catalog uniformly and therefore still pollute the cache.

\subsection{Traffic Mixing and Ratios}
To evaluate the DTAC system under varying levels of infrastructure strain, the generator blends the human sessions and the bot traffic into chronological, single-stream request logs. For a target bot ratio $r$ and $H \approx 100{,}000$ human measurement requests, a file contains $B = H \cdot r / (1 - r)$ bot requests spread evenly over the entire measurement window, so that the bot share is approximately $r$ throughout the timeline rather than concentrated at the start.

The traffic was multiplexed to create multiple distinct datasets, each representing a different severity of AI interference. The test ratios were defined by the percentage of total traffic originating from simulated scrapers:
\begin{itemize}
    \item \textbf{10\% Bot Traffic:} Representing a baseline, low-threat environment.
    \item \textbf{30\% Bot Traffic:} Representing moderate, localized scraper activity.
    \item \textbf{50\% Bot Traffic:} Representing a heavy data-harvesting event.
    \item \textbf{75\% Bot Traffic:} Representing a severe scan event capable of inducing cache thrashing \cite{zhang_rethinking_2026}.
\end{itemize}

Each log is prefixed with a warm-up phase of about 50,000 purely human requests (Section 3.6.1). For every seed, this yields eight files (two bot workloads at four ratios). Alongside each file the generator writes a metadata file recording the seed, every generator parameter and the generator's version, which the replayer copies into its results.

\section{Experimental Setup and Benchmarking Procedures}
To empirically validate the efficacy of the DTAC algorithm against standard and state-of-the-art caching mechanisms, a controlled benchmarking protocol was established and executed on the laptop described in Section 3.3.1.

\subsection{The Benchmarking Protocol}
The lifecycle of each individual experiment followed a three-step protocol to ensure that the collected data reflected steady-state behavior rather than artificial testing anomalies.

\begin{enumerate}
    \item \textbf{Step 1: Cache Warm-Up:} A cache starting from an empty state (a ``cold'' cache) naturally experiences an initial 100\% miss rate as it begins storing objects, which skews performance metrics. Therefore, before any simulated scraper traffic, each log begins with a fixed prefix of about 50,000 purely human requests. These requests are processed by the cache exactly like any other request, so the cache is populated, but they are excluded from every reported metric and timing.
    \item \textbf{Step 2: Mixed Traffic Replay:} Immediately after the warm-up, the replayer continues with the measurement phase, in which human and bot requests are interleaved at the file's bot ratio. Each combination of seed, workload, bot ratio, algorithm and capacity was run as a separate replay, and the replays were executed sequentially so that runs never competed for CPU.
    \item \textbf{Step 3: Metric Logging:} Throughout the measurement phase, the replayer recorded the outcome of every request (hit or miss, bytes served, latency of each phase, and the router's score), sampled the contents of the cache periodically to measure pollution, and recorded every main-cache eviction together with the label (human or bot) of the request that had caused the evicted entry to be admitted. Each run wrote all of its metrics, its configuration and the software version to a JSON file.
\end{enumerate}

The experiment grid comprised five seeds (1--5), two bot workloads (naive and stealthy), four bot ratios (10\%, 30\%, 50\% and 75\%) and five algorithms at a primary cache capacity of 10\,000 entries (10\% of the catalog). In addition, a capacity sweep at 1\,000, 5\,000, 10\,000 and 20\,000 entries (1--20\% of the catalog) was run for the naive workload at 30\% and 75\% bot traffic. In total the grid contained 350 replays. A single replay was also repeated over the network (Section 3.3.3) to validate the in-process results.

\subsection{Baseline Comparison}
To isolate and quantify the performance impact of the AI-aware design, every workload was replayed through five algorithms:

\begin{itemize}
    \item \textbf{Standard LRU (Control):} A sharded Least Recently Used cache that ignores the AI-score entirely and treats all incoming traffic (both human and bot) with equal temporal importance. This establishes how severely a standard cache degrades under scan patterns.
    \item \textbf{SIEVE and S3-FIFO (State-of-the-art baselines):} The two modern scan-resistant algorithms reviewed in Chapter 2 \cite{zhang_sieve_2024, yang_fifo_2023}, implemented for this study. They are also blind to the AI-score and infer the value of objects from access patterns alone.
    \item \textbf{DTAC without the score (Ablation):} The DTAC engine with every request assigned the same score a clean human request would receive ($\text{sigmoid}(L_0) = 0.02$), so bots cannot be told apart from humans. Using this constant, rather than zero, keeps the promotion speed identical to DTAC's for human traffic. Comparing DTAC with this ablation separates the contribution of the AI-score from that of the ghost cache and the decayed-score tournament.
    \item \textbf{DTAC (Experimental):} The full system, in which the router's AI-score weights every request's contribution to admission and eviction.
\end{itemize}

All five algorithms share the same shard count and the same per-shard capacity split, so only the caching policy differs between them. Because the request sequence, the total cache capacity and the hardware were identical across the five runs of a given cell, any divergence in the final performance metrics can be attributed to the algorithms themselves.

\subsection{Statistical Analysis}
Each metric is reported as the mean over the five evaluation seeds with a 95\% confidence interval computed from the $t$-distribution ($n = 5$). Because every algorithm processes exactly the same human requests for a given seed, comparisons between DTAC and each alternative are made as \emph{paired differences}: the difference $\text{metric}(\text{DTAC}) - \text{metric}(X)$ is computed separately for each seed, and the mean and 95\% confidence interval of these differences are reported. A difference whose confidence interval excludes zero is treated as statistically significant; one whose interval includes zero is reported as a tie. This paired design is considerably more sensitive than comparing two overlapping confidence intervals.

\section{Evaluation Metrics}
To evaluate the Dynamic Traffic-Aware Caching (DTAC) algorithm against the baselines, the benchmarking protocol recorded a specific set of quantitative metrics. These metrics were designed to assess both the protective capability of the cache in preserving human user experience and the computational overhead introduced by the heuristic evaluations.

\subsection{Cache Efficiency Metrics}
The primary goal of the DTAC system is not necessarily to increase the overall cache hit rate, but to preserve the hit rate for legitimate human users during an active AI scraper attack. Therefore, the hit rates were isolated by traffic origin using the ground-truth labels in the logs.

\textbf{Human Cache Hit Rate ($CHR_{human}$):} This was the primary metric of the study. It measures the percentage of legitimate human requests successfully served from the cache rather than requiring a costly fetch from the origin server. It is calculated as:

\begin{equation}
CHR_{human} = \frac{H_{hit}}{H_{total}}
\end{equation}

Where $H_{hit}$ is the number of human requests found in the cache, and $H_{total}$ is the total number of human requests in the measurement phase. A successful DTAC implementation should maintain a stable $CHR_{human}$ as the volume of bot traffic increases, at least when the bots are detectable. The overall hit rate (over all requests) was recorded as a secondary metric, and the bot hit rate for reference only: a low bot hit rate is not a goal in itself, provided it does not come at the expense of the human hit rate.

\textbf{Byte Hit Rate (BHR):} While object hit rate measures individual files, BHR measures the total volume of data served from the cache. This is a relevant metric for evaluating bandwidth and egress cost savings. It is defined as:

\begin{equation}
BHR = \frac{\sum B_{hit}}{\sum B_{total}}
\end{equation}

Where $B_{hit}$ represents the bytes served from the cache and $B_{total}$ represents the total bytes requested. BHR was computed for human requests and overall. Because cache capacity in this study is measured in entries rather than bytes, and none of the evaluated algorithms takes object size into account, BHR reflects how hits on popular items happen to correlate with object size rather than the behavior of a size-aware policy.

\subsection{Cache Pollution Index (CPI)}
To quantify how much cache capacity the simulated scrapers wasted, a Cache Pollution Index (CPI) was tracked. The CPI measures the proportion of resident cache entries that no human has requested since they were admitted:

\begin{equation}
CPI = \frac{M_{bot}}{M_{total}}
\end{equation}

Where $M_{bot}$ is the number of resident main-cache entries that have not been requested by any human since entering the cache, and $M_{total}$ is the number of resident main-cache entries. The CPI was sampled every 10,000 measurement requests and at the end of the replay. This definition deliberately does not count an entry as pollution merely because a bot request happened to trigger its admission: if humans subsequently use the entry, the capacity is not wasted. The simpler measure, the share of resident entries whose admission was triggered by a bot request, was recorded alongside it as the \emph{bot-triggered share}.

For the standard LRU baseline, the CPI is expected to grow with the injected bot traffic ratio, since every bot request is admitted. For DTAC, content requested only by detected bots should rarely accumulate enough score to be promoted, so its CPI should remain close to zero when the router identifies the bots correctly.

\subsection{Router Classification Quality}
Because the benefit of DTAC depends on the quality of the AI-score, the router was also evaluated as a binary classifier against the ground-truth labels, independently of the cache. The metrics recorded were the area under the Receiver Operating Characteristic curve (ROC-AUC) of the bot probability; precision, recall (true-positive rate) and false-positive rate at a decision threshold of 0.5; the false-positive rate separately for new and returning human visitors; the mean bot probability of each class and the gap between the two means, which directly determines how differently the cache treats the two classes; and, for each heuristic, how often it fired on human and on bot requests.

\subsection{Computational Latency and Overhead}
Because the proxy and cache sit inline with the network traffic, any algorithmic modifications must not introduce significant delays. While Rust's performance characteristics keep the underlying execution fast, the algorithmic overhead still required explicit measurement.

\begin{itemize}
    \item \textbf{Proxy Scoring Overhead ($T_{proxy}$):} The time taken for the Heuristic Proxy to update the client's state and calculate the AI-score. Since this relies on bounded $\mathcal{O}(\log K)$ data structures, $T_{proxy}$ should remain well below 2 milliseconds.
    \item \textbf{Cache Overhead ($T_{cache}$):} The time required for the cache engine to look up the resource, update its score, and perform any ghost-cache bookkeeping, promotion and eviction, including the (instantaneous, in the in-process replay) origin fetch.
    \item \textbf{Total Access Time:} The end-to-end time of a cache access, recorded for every algorithm.
\end{itemize}

These durations were measured with Rust's high-precision \texttt{std::time::Instant} and recorded in nanoseconds into \texttt{hdrhistogram} histograms, from which the 50th, 90th and 99th percentiles and the maximum were reported. Memory overhead was measured as the growth of the process's resident set size (RSS) over the replay, together with the number of main-cache and ghost-cache entries; because all cached payloads share a single buffer in the replayer, the RSS growth reflects cache metadata rather than payload bytes.
